\subsection{Large Extra Dimensions and neutrino experiments --- arXiv:2207.02790}

\textbf{Authors:} D. V. Forero, C. Giunti, C. A. Ternes, O. Tyagi

\paragraph{主な主張}
最小LED模型の結合解析から$R<0.20\,\mu\mathrm{m}$（正常順序）および$R<0.10\,\mu\mathrm{m}$（逆順序）を90\%信頼水準で得る一方、Gallium異常が選ぶ領域は他の振動データと強く緊張することを示す\cite{Forero:2022LEDExperiments}。

\paragraph{新規性}
BESTを含むGalliumデータ、複数の原子炉フラックス模型、MINOS/MINOS+・Daya Bay・KATRINを比較し、異常の説明と他実験による制約を同じLEDパラメータ空間で再評価した。

\paragraph{位置付け}
最小LEDに対する2022年時点の包括的な実験制約の基準であり、半径制約を引用する際の模型仮定と質量順序を明示するために重要である。

\paragraph{KATRIN解析}
第2回KATRIN測定のトリチウム$\beta$崩壊の終端から40\,eVの積分電子スペクトルを用い、娘分子の終状態分布を含むKurie関数に、各KKモードの質量$m_i^{(n)}$と混合重み$|U_{ei}|^2(V_{in})^2$を組み込んで解析した。質量40\,eV未満のKKモードを足し合わせ、規格化や背景などの系統不確かさを考慮したフィットから、$(R_{\mathrm{ED}},m_0)$平面に独立の制限を与えている\cite{Forero:2022LEDExperiments}。
